% © 2026 Ing. David Jaroš — CC BY-NC-ND 4.0
%
% This work is licensed under a Creative Commons Attribution-NonCommercial-NoDerivatives
% 4.0 International License (CC BY-NC-ND 4.0).
%
% License History: Earlier drafts (up to v0.3) were released under CC BY 4.0.
% From v0.4 onward, all material is released under CC BY-NC-ND 4.0 to protect
% the integrity of the theoretical work during ongoing academic development.
%
% See LICENSE.md for full license text.
%
% Track: T1_GR — General Relativity Recovery
% Purpose: Paper-ready theorem chain Θ → g → Γ → R → Einstein
% Sources: canonical/gr_closure/, canonical/t_munu/, canonical/geometry/

% UBT-AI-PROVENANCE-BEGIN
% schema: ubt-ai-provenance/v1
% tier: C_working
% ai_assistance: disclosed
% human_review: risk-based
% editorial_responsibility: Ing. David Jaroš
% policy: ../../AI_PROVENANCE.md
% notice: Working material; exhaustive human review is not claimed.
% UBT-AI-PROVENANCE-END

\documentclass[12pt,a4paper]{article}
\usepackage[a4paper,margin=2.5cm]{geometry}
\usepackage{amsmath,amssymb,amsthm}
\usepackage{hyperref}
\usepackage{tcolorbox}
\tcbuselibrary{skins,breakable}
\usepackage{booktabs}

\newtheorem{theorem}{Theorem}[section]
\newtheorem{lemma}[theorem]{Lemma}
\newtheorem{corollary}[theorem]{Corollary}
\newtheorem{proposition}[theorem]{Proposition}
\theoremstyle{definition}
\newtheorem{definition}[theorem]{Definition}
\theoremstyle{remark}
\newtheorem{remark}[theorem]{Remark}

\newcommand{\Lzero}{\textbf{[L0]}}
\newcommand{\Lone}{\textbf{[L1]}}
\newcommand{\Ltwo}{\textbf{[L2]}}
\newcommand{\OPEN}{\textbf{[OPEN]}}

\title{\textbf{GR Recovery in UBT: Complete Theorem Chain}\\
  \large $\Theta \to g_{\mu\nu} \to \Gamma^\lambda_{\mu\nu}
         \to R^\rho{}_{\sigma\mu\nu} \to G_{\mu\nu} = 8\pi G\,T_{\mu\nu}$}
\author{Ing.\ David Jaroš}
\date{April 2026}

\begin{document}
\maketitle
\tableofcontents

% ============================================================
\section{Foundations: UBT Axioms and Field}
\label{sec:foundations}
% ============================================================

\begin{definition}[Biquaternion algebra]
\label{def:biquaternion_algebra}
The biquaternion algebra is $\mathbb{B} := \mathbb{C}\otimes_{\mathbb{R}}\mathbb{H}$.
As a vector space over $\mathbb{R}$, $\dim_{\mathbb{R}}\mathbb{B} = 8$.
There is a canonical algebra isomorphism
\[
  \mathbb{C}\otimes_{\mathbb{R}}\mathbb{H} \;\cong\; \mathrm{Mat}(2,\mathbb{C}).
\]
Ref: \texttt{canonical/fields/biquaternion\_algebra.tex} \Lzero.
\end{definition}

\begin{definition}[Complex time and AXIOM B]
\label{def:complex_time}
Complex time is $\tau := t + i\psi \in \mathbb{C}$, where $t \in \mathbb{R}$ is
real time and $\psi \in \mathbb{R}$ is the imaginary (phase) component.
\textbf{AXIOM B}: the complex-time derivative $\partial_\tau$ lies in the timelike
sector of $\mathrm{Cl}_{1,3}(\mathbb{R})$, so that
\[
  \langle\partial_\tau, \partial_\tau\rangle_{\eta} < 0.
\]
Ref: \texttt{canonical/THEORY/axioms/core\_assumptions.tex} \Lzero.
\end{definition}

\begin{definition}[Fundamental field and admissible class]
\label{def:theta_field}
The fundamental UBT field is a map
\[
  \Theta : M^4 \times \mathbb{C}_\tau \;\longrightarrow\; \mathbb{B},
\]
satisfying the T-shirt equation
\[
  \nabla^\dagger\nabla\,\Theta(q,\tau) \;=\; \kappa\,\mathcal{T}(q,\tau).
\]
The admissible field class is
\[
  \mathcal{A}_{\mathrm{UBT}}
  := \bigl\{\Theta \;\big|\; \partial_\mu\Theta \text{ linearly independent},\;
     \Theta \neq \mathrm{const}\bigr\}.
\]
Ref: \texttt{canonical/fields/theta\_field.tex} \Lzero.
\end{definition}

% ============================================================
\section{Step 1: Metric Emergence}
\label{sec:step1}
% ============================================================

\begin{theorem}[Metric emergence, \Lone]
\label{thm:metric_emergence}
For $\Theta \in \mathcal{A}_{\mathrm{UBT}}$, define the biquaternionic metric
tensor
\[
  \mathcal{G}_{\mu\nu} \;:=\; \partial_\mu\Theta \cdot \partial_\nu\Theta^\dagger,
\]
and the derived real metric
\[
  g_{\mu\nu} \;:=\; \frac{\mathrm{Re}\bigl[\mathrm{Tr}(\partial_\mu\Theta
  \cdot\partial_\nu\Theta^\dagger)\bigr]}{\mathcal{N}},
  \qquad \mathcal{N} := \mathrm{Re}\bigl[\mathrm{Tr}
  (\partial_0\Theta\cdot\partial_0\Theta^\dagger)\bigr] > 0.
\]
The tensor $g_{\mu\nu}$ is symmetric and transforms as a covariant rank-2 tensor
under coordinate changes on $M^4$.

\medskip
\emph{Proof sketch}: Symmetry of $g_{\mu\nu}$: $\mathrm{Re}[\mathrm{Tr}(AB^\dagger)]
= \mathrm{Re}[\mathrm{Tr}(BA^\dagger)]$ for $A,B\in\mathrm{Mat}(2,\mathbb{C})$.
Covariance: $\partial_\mu\Theta$ transforms as a covariant vector under
diffeomorphisms.  See \texttt{canonical/gr\_closure/step1\_metric\_bridge.tex}.
\hfill$\square$
\end{theorem}

% ============================================================
\section{Step 2: Non-Degeneracy}
\label{sec:step2}
% ============================================================

\begin{theorem}[Non-degeneracy of $g_{\mu\nu}$, \Lone]
\label{thm:nondegeneracy}
For $\Theta \in \mathcal{A}_{\mathrm{UBT}}$,
\[
  \det(g_{\mu\nu}) \;\neq\; 0.
\]
Equivalently, the Gram matrix $G_{\mu\nu} := \mathrm{Re}[\mathrm{Tr}(\partial_\mu\Theta
\cdot\partial_\nu\Theta^\dagger)]$ is invertible if and only if the four derivatives
$\partial_\mu\Theta$ are linearly independent in $\mathbb{B}$.

\medskip
\emph{Proof sketch}: The Gram matrix of a set of $n$ vectors $v_\mu$ over an inner
product space satisfies $\det(G) = 0$ iff the vectors are linearly dependent.
Linear independence is the defining condition of $\mathcal{A}_{\mathrm{UBT}}$.
See \texttt{canonical/gr\_closure/step2\_nondegeneracy.tex}, Theorem~1.
\hfill$\square$
\end{theorem}

% ============================================================
\section{Step 3: Lorentzian Signature}
\label{sec:step3}
% ============================================================

\begin{theorem}[Lorentzian signature from AXIOM B, \Lone]
\label{thm:signature}
The derived metric $g_{\mu\nu}$ has Lorentzian signature $(-,+,+,+)$.
Specifically: $g_{00} < 0$ and the spatial block $(g_{ij})_{i,j=1,2,3}$
is positive-definite.

\medskip
\emph{Proof sketch}: By AXIOM~B (Definition~\ref{def:complex_time}),
$\partial_\tau$ lies in the timelike sector of $\mathrm{Cl}_{1,3}(\mathbb{R})$.
Under the identification $\tau = t + i\psi$, the temporal derivative
$\partial_t\Theta = \partial_t(\mathrm{Re}\,\Theta) + i\partial_t(\mathrm{Im}\,\Theta)$
inherits the Clifford-algebraic timelike property, giving
$g_{00} = \mathrm{Re}[\mathrm{Tr}(\partial_t\Theta\cdot\partial_t\Theta^\dagger)]/\mathcal{N} < 0$.
The spatial components satisfy $g_{ii} > 0$ by the spacelike nature of the
spatial Clifford generators.
The normalization $\mathcal{N}$ does \textbf{not} determine the signature;
signature is an algebraic theorem.
See \texttt{canonical/gr\_closure/step3\_signature\_theorem.tex}, Theorem~2.
\hfill$\square$
\end{theorem}

\begin{remark}
The Lorentzian signature is a \emph{theorem}, not a postulate.  It follows
from AXIOM~B alone, independent of the specific form of $\Theta$.
\end{remark}

% ============================================================
\section{Step 4: Standard GR Geometric Apparatus}
\label{sec:step4}
% ============================================================

\begin{proposition}[Levi-Civita connection in the real sector after projection, standard]
\label{prop:levi_civita}
Given the non-degenerate Lorentzian metric $g_{\mu\nu}$ from Theorem~\ref{thm:metric_emergence},
the unique torsion-free metric-compatible connection is the Levi-Civita connection in the real projection
\[
  \Gamma^\lambda_{\mu\nu} \;=\; \tfrac{1}{2}g^{\lambda\rho}
  \bigl(\partial_\mu g_{\nu\rho} + \partial_\nu g_{\mu\rho} - \partial_\rho g_{\mu\nu}\bigr).
\]
\end{proposition}

\begin{proposition}[Curvature tensors, standard]
\label{prop:curvature}
The Riemann tensor, Ricci tensor, Ricci scalar, and Einstein tensor are defined from
$\Gamma^\lambda_{\mu\nu}$ by the standard formulas of differential geometry:
\begin{align}
  R^\rho{}_{\sigma\mu\nu} &:= \partial_\mu\Gamma^\rho_{\nu\sigma}
    - \partial_\nu\Gamma^\rho_{\mu\sigma}
    + \Gamma^\rho_{\mu\lambda}\Gamma^\lambda_{\nu\sigma}
    - \Gamma^\rho_{\nu\lambda}\Gamma^\lambda_{\mu\sigma}, \\
  R_{\mu\nu} &:= R^\rho{}_{\mu\rho\nu}, \qquad R := g^{\mu\nu}R_{\mu\nu}, \\
  G_{\mu\nu} &:= R_{\mu\nu} - \tfrac{1}{2}g_{\mu\nu}R.
\end{align}
The Bianchi identity $\nabla^\mu G_{\mu\nu} = 0$ holds by standard geometry.
\end{proposition}

\begin{remark}
In UBT these objects are all \emph{derived} (not postulated) quantities,
obtained as real projections:
$g_{\mu\nu} = \mathrm{Re}(\mathcal{G}_{\mu\nu})$,
$\Gamma = \mathrm{Re}(\Omega)$, etc.
Ref: \texttt{canonical/geometry/curvature.tex}.
\end{remark}

% ============================================================
\section{Step 5: Einstein Field Equations}
\label{sec:step5}
% ============================================================

\begin{theorem}[Einstein field equations, \Lone]
\label{thm:einstein}
Consider the total UBT action
\[
  S_{\mathrm{total}}[g,\Theta]
  \;=\; \frac{1}{16\pi G}\int \sqrt{-g}\,R\,\mathrm{d}^4x
  \;+\; S_\Theta[g,\Theta],
\]
where $S_\Theta$ is the matter action for $\Theta$ with kinetic term
$\mathrm{Tr}[(D_\mu\Theta)^\dagger D^\mu\Theta]$.
Variation with respect to $g^{\mu\nu}$ gives
\[
  G_{\mu\nu} \;=\; 8\pi G\,T_{\mu\nu},
\]
where the stress-energy tensor is
\[
  T_{\mu\nu} \;=\; \mathrm{Re}\bigl(\partial_\mu\Theta\,\partial_\nu\Theta^\dagger\bigr)
  - \tfrac{1}{2}g_{\mu\nu}\,g^{\alpha\beta}
    \mathrm{Re}\bigl(\partial_\alpha\Theta\,\partial_\beta\Theta^\dagger\bigr),
\]
satisfying $T_{\mu\nu} = T_{\nu\mu}$ and $\nabla^\mu T_{\mu\nu} = 0$.

\medskip
\emph{Proof sketch}: Standard Hilbert variation.  See
\texttt{canonical/t\_munu/step3\_einstein\_with\_matter.tex}.
The $T_{\mu\nu}$ symmetry and conservation are proved in
\texttt{canonical/geometry/stress\_energy.tex}.
\hfill$\square$
\end{theorem}

% ============================================================
\section{Schwarzschild Metric from $\Theta_0$}
\label{sec:schwarzschild}
% ============================================================

\begin{theorem}[Schwarzschild metric, \Lone]
\label{thm:schwarzschild}
The ansatz
\[
  \Theta_0 \;=\; e^{i\Phi(r)}\bigl[f(r)\,\mathbf{1} + g(r)\,\boldsymbol{e}_r\bigr],
\]
with $\tau = t + i\psi$, $g(r) = r\Psi(r)^2$, $f'(r) = \Psi(r)\sqrt{2M/r}$,
and $\Phi(r) = \alpha(r) = (1-M/2r)/(1+M/2r)$, reproduces the Schwarzschild
metric in isotropic coordinates:
\[
  g_{tt} = -\Phi(r)^2, \qquad g_{ij} = \Psi(r)^4\,\delta_{ij},
  \qquad \Psi(r) = 1 + \frac{M}{2r}.
\]
\emph{Numerical verification}: \texttt{tools/verify\_schwarzschild\_theta.py},
relative error $< 10^{-8}$.
Ref: \texttt{canonical/geometry/biquaternionic\_vacuum\_solutions.tex~\S3}.
\end{theorem}

\begin{theorem}[ASD condition and twistor space, \Lone]
\label{thm:asd}
For $\Theta \in \mathrm{SU}(2)_- \subset \mathbb{C}\otimes\mathbb{H}$ smooth with
$|\Theta|=1$:
\begin{enumerate}
  \item The holonomy of $g_{\mu\nu}[\Theta]$ lies in $\mathrm{Sp}(1)\cong\mathrm{SU}(2)_-$,
        which implies the anti-self-dual (ASD) Weyl condition $C^+ = 0$.
  \item Combined with the vacuum UBT equation $\nabla^\dagger\nabla\Theta = 0$,
        giving $R_{\mu\nu} = 0$ via the GR chain, the metric is ASD Ricci-flat.
  \item By the Penrose nonlinear graviton theorem, $g_{\mu\nu}[\Theta]$ admits a
        curved twistor space description.
\end{enumerate}
Note: Schwarzschild (Petrov type D) lies outside the $\mathrm{SU}(2)_-$ sector.
\end{theorem}

% ============================================================
\section{Regge-Wheeler Equation (Linearised Gravity)}
\label{sec:regge_wheeler}
% ============================================================

\begin{theorem}[Linearised GR and Regge-Wheeler, \Lone]
\label{thm:regge_wheeler}
Linearisation of the UBT field equation about flat background
$\Theta = \Theta_0 + \epsilon\,\delta\Theta$ reproduces the linearised
Einstein equations.  For perturbations of Schwarzschild in the odd-parity
(axial) sector with angular decomposition, the perturbation equation reduces
to the Regge-Wheeler equation
\[
  \left[\frac{\mathrm{d}^2}{\mathrm{d}r_*^2} + \omega^2
        - V_{\mathrm{RW}}(r)\right]\Psi_{\mathrm{RW}} = 0,
\]
where $r_*$ is the tortoise coordinate, $\omega$ the graviton frequency, and
$V_{\mathrm{RW}}$ the Regge-Wheeler potential.
This result follows without additional input from UBT beyond the metric derivation.
\end{theorem}

\begin{tcolorbox}[colback=orange!5!white,colframe=orange!60!black,
  title=Open: Zerilli equation (even-parity) \OPEN~\Ltwo]
The even-parity (polar) perturbation equation of Schwarzschild — the Zerilli
equation — has \textbf{not} yet been derived from UBT.
This is documented as an open problem [L2] and does not affect the main GR
recovery result.  A complete treatment would require extending the analysis of
even-parity metric perturbations in the biquaternion framework.
\end{tcolorbox}

% ============================================================
\section{Open Problem: Off-Shell $\Theta$-Only Closure}
\label{sec:gap}
% ============================================================

\begin{tcolorbox}[colback=red!5!white,colframe=red!60!black,
  title=Open Problem GAP-10: Off-shell $\Theta$-only closure \OPEN~\Ltwo]

\textbf{On-shell result (proved)}: For $\Theta \in \mathcal{A}_{\mathrm{UBT}}$
satisfying its own Euler--Lagrange equation, and assuming the gauge-reduced
local rank condition, $\delta\hat{S}/\delta\Theta = 0$ is equivalent to the
Einstein equations evaluated on $g = g[\Theta]$.

\medskip
\textbf{Off-shell gap}: Global non-degeneracy of $J = \delta g^{\mu\nu}/\delta\Theta$
for \emph{all} field configurations (not only on-shell solutions in
$\mathcal{A}_{\mathrm{UBT}}$) is \textbf{not proved}.

\medskip
\textbf{Missing lemma}: Show that $\ker J$ consists only of gauge directions
(pure phase or diffeomorphism) for all $\Theta$ in the full off-shell field space.

\medskip
\textbf{Obstructions identified}:
\begin{itemize}
  \item Rank mismatch: $\mathrm{Re}(\nabla^\dagger\nabla\Theta)$ is rank-0;
        $G_{\mu\nu}$ is rank-2 — direct identification fails.
  \item Topology: global injectivity of $\Theta\to g[\Theta]$ requires
        $H^2(M^4,\mathbb{Z})$ analysis of the $\Theta$-bundle.
  \item Non-perturbative: fixed-point theorem needed in appropriate Sobolev space.
\end{itemize}

\medskip
Ref: \texttt{research\_tracks/research/gr\_offshell\_gap.md},
\texttt{canonical/gr\_closure/step2\_theta\_only\_closure.tex}.
This gap does \textbf{not} affect the validity of Theorems~\ref{thm:metric_emergence}--\ref{thm:regge_wheeler}.
\end{tcolorbox}

% ============================================================
\section{Summary Table}
\label{sec:summary}
% ============================================================

\begin{tcolorbox}[colback=blue!3!white,colframe=blue!60!black,title=GR Chain: Complete Status]
\begin{center}
\begin{tabular}{clll}
\toprule
\textbf{Step} & \textbf{Claim} & \textbf{Status} & \textbf{Reference} \\
\midrule
1 & Metric $g_{\mu\nu}$ from $\Theta$ & Proved \Lone & step1\_metric\_bridge.tex \\
2 & Non-degeneracy $\det(g)\neq 0$ & Proved \Lone & step2\_nondegeneracy.tex \\
3 & Signature $(-,+,+,+)$ from AXIOM B & Proved \Lone & step3\_signature\_theorem.tex \\
4 & $g\to\Gamma\to R$ geometric chain & Standard & Diff.\ geometry \\
5 & Einstein eqs $G_{\mu\nu}=8\pi GT_{\mu\nu}$ & Proved \Lone & step3\_einstein\_with\_matter.tex \\
6a & Schwarzschild metric & Proved \Lone & biquaternionic\_vacuum\_solutions.tex \\
6b & ASD condition / twistor & Proved \Lone & asd\_condition\_ubt.tex \\
6c & Regge-Wheeler equation & Proved \Lone & (linearized GR) \\
7a & Zerilli equation (even-parity) & Open \Ltwo & — \\
7b & Off-shell $\Theta$-only closure & Open \Ltwo & gr\_offshell\_gap.md \\
\bottomrule
\end{tabular}
\end{center}
\end{tcolorbox}

\end{document}
